\chapter{Implementación}

En este capítulo se describe la implementación del sistema desarrollado, abordando tanto las decisiones técnicas adoptadas durante la codificación como los aspectos más relevantes o complejos del proceso. La exposición se organiza en torno a tres ejes principales: la estructura del código y la organización del proyecto, los detalles de implementación más significativos desde el punto de vista técnico, y el sistema de bot, que constituye la parte técnicamente más exigente del trabajo. Finalmente, se recogen las licencias de todos los recursos externos utilizados.

% ─────────────────────────────────────────────────────────────────────────────
\section{Organización del código}
% ─────────────────────────────────────────────────────────────────────────────

Para garantizar la escalabilidad del proyecto y facilitar el mantenimiento del código, se ha seguido una estructura de directorios basada en la separación de responsabilidades. Dado que Godot utiliza un sistema de escenas y nodos, la organización se ha diseñado para diferenciar claramente entre los elementos visuales, la lógica del videojuego pura y la infraestructura de datos.

A continuación, se detallan las carpetas principales que componen el proyecto y su propósito dentro de la arquitectura global:

\begin{itemize}
	\item \textbf{Assets/}: Esta carpeta centraliza todos los recursos gráficos y visuales del videojuego. Contiene las texturas (sprites), fuentes tipográficas, temas de interfaz (Themes) y recursos específicos de Godot como los \texttt{TileMaps}. Al agrupar estos elementos, se consigue que cualquier cambio estético o de diseño visual sea fácil de localizar sin interferir con la lógica programática.
	
	\item \textbf{BotEnvironment/}: Contiene todo el entorno necesario para el entrenamiento y la ejecución del agente de Inteligencia Artificial. Aquí se incluyen los scripts y configuraciones que permiten al \textit{bot} interactuar con el videojuego, separando el comportamiento del agente del resto de sistemas para permitir pruebas aisladas de la IA.
	
	\item \textbf{Data \& Infrastructure/}: En este directorio se ubican la definición de estructuras de datos y servicios de infraestructura que no pertenecen a la lógica del videojuego, pero que son fundamentales para el funcionamiento del sistema, como el manejo del audio y la navegación entre escenas.
	
	\item \textbf{UI/}: Siguiendo el patrón de diseño orientado a escenas de Godot, esta carpeta contiene tanto los scripts de control de la interfaz de usuario como las escenas (.tscn) correspondientes a los menús, paneles y elementos interactivos (HUD). Esto permite que el diseño de la interfaz esté estrechamente ligado a su lógica de control inmediata, pero separado de la lógica de combate o movimiento.
	
	\item \textbf{GameLogic/}: Es el núcleo del videojuego y presenta una subdivisión estratégica para cumplir con el diseño arquitectónico:
	\begin{itemize}
		\item \textbf{Lógica Pura (C\# scripts):} Contiene clases puras de C\# que no heredan de la clase \texttt{Node} de Godot. Aquí residen las reglas del videojuego y algoritmos que podrían ser testeados de forma unitaria fuera del motor.
		\item \textbf{Nodos de Godot:} Subdividida en carpetas como Pieces, Tiles y Board. En cada una de ellas se encuentran las escenas de Godot y los scripts que heredan de la API del motor para gestionar la representación espacial y el comportamiento de los objetos en pantalla.
	\end{itemize}
\end{itemize}

Esta estructura modular permite que el proyecto sea robusto, ya que cada componente tiene una responsabilidad bien definida, permitiendo que el motor de videojuegos funcione de manera fluida y organizada.

% ─────────────────────────────────────────────────────────────────────────────
\section{Licencias requeridas}
% ─────────────────────────────────────────────────────────────────────────────

Para el desarrollo de \textit{NeuroForge}, se han utilizado diversos recursos externos de audio y gráficos, así como el propio motor de desarrollo, bajo licencias que permiten su uso en proyectos académicos y comerciales. A continuación, se detallan las licencias de software y recursos empleadas.

\subsection{Motor de Videojuegos: Godot Engine}

El proyecto ha sido desarrollado utilizando \textit{Godot Engine}. Este motor se distribuye bajo la \textit{Licencia MIT}, una de las licencias de software libre más permisivas y utilizadas en la industria.

\begin{itemize}
	\item \textbf{Permisos:} Permite el uso, copia, modificación, fusión, publicación, distribución, sublicencia y venta del software sin restricciones.
	\item \textbf{Condiciones:} El único requisito es incluir el aviso de copyright y el texto de la licencia original de Godot en cualquier copia sustancial del software o en la documentación del mismo.
\end{itemize}

\subsection{Recursos de Audio}

Los efectos de sonido y la música han sido obtenidos respetando las licencias especificadas por los autores en plataformas de contenido libre:

\begin{itemize}
	\item \textbf{Música durante la partida:} \textit{Granular 5.wav} por PodcastAC. Disponible en Freesound\footnote{\url{https://freesound.org/s/663786/}}. Licencia: \textit{Attribution 4.0 (CC BY 4.0)}.
	\item \textbf{Música del menú:} \textit{Ambient, Piano, Atmosphere - Reborn} por Andrewkn. Disponible en Freesound\footnote{\url{https://freesound.org/s/720772/}}. Licencia: \textit{Attribution 4.0 (CC BY 4.0)}.
	\item \textbf{Efectos de sonido de las piezas:} \textit{SCIMisc\_Energy UI Low 01} por KVV\_Audio. Disponible en Freesound\footnote{\url{https://freesound.org/s/811930/}}. Licencia: \textit{Attribution 4.0 (CC BY 4.0)}.
	\item \textbf{Interfaz de usuario:} \textit{UIClick\_Button press} por newlocknew. Disponible en Freesound\footnote{\url{https://freesound.org/s/665222/}}. Licencia: \textit{Creative Commons 0 (CC0)}.
	\item \textbf{Interacción con el tablero:} Recursos obtenidos de \textit{Pixabay} bajo su \textit{Content License}, que permite el uso gratuito y modificación sin atribución obligatoria.
\end{itemize}

\subsection{Recursos Gráficos}

El apartado visual integra \textit{assets} adaptados y modificados bajo los siguientes términos:

\begin{itemize}
	\item \textbf{CraftPix Assets:} Los recursos provenientes de \textit{Craftpix.net}\footnote{\url{https://craftpix.net/sets/cyberpunk-platformer-asset-pixel-art/}} se utilizan bajo su licencia comercial, que permite su uso ilimitado en videojuego siempre que no se revendan los archivos fuente originales. Se han utilizado para la mayoria de los recursos gráficos dentro del videojuego.
	\item \textbf{Sprites de piezas (Mechs):} \textit{Mech Assets} por \textit{Marinesosa}. Disponible en itch.io\footnote{\url{https://marinesosa.itch.io/mech-assets}}. Su licencia permite la modificación y el uso comercial/no comercial, prohibiendo la reventa del paquete original. Se han utilizado para el diseño de algunas piezas
\end{itemize}

\subsection{Modificación de recursos}

Es importante señalar que gran parte de los recursos gráficos y de audio han sido modificados para ajustarse a las necesidades de \textit{NeuroForge}. Estas modificaciones (ajustes de color, reescalado, edición de audio) se han realizado mediante \textit{Aseprite} y \textit{Audacity}, cumpliendo con los permisos de obras derivadas de las licencias mencionadas.

% ─────────────────────────────────────────────────────────────────────────
\section{Detalles de implementación relevantes}
% ─────────────────────────────────────────────────────────────────────────

Entre todos los sistemas implementados, la arquitectura del bot constituye el componente técnicamente más exigente del proyecto, tanto por la complejidad del entorno de entrenamiento como por su integración con el motor de videojuegos.

\subsection{Arquitectura del bot}

El bot presenta dos modos de funcionamiento. El comportamiento por defecto, utilizado como \textit{fallback} cuando el modelo entrenado no está disponible o falla al cargar, consiste en la selección aleatoria de un movimiento válido entre todos los posibles en cada turno. Este modo garantiza que la partida siempre pueda completarse y sirve como línea base para evaluar la mejora que aporta el modelo de inteligencia artificial. Esto fue sencillo de hacer ya que además se utilizo el modelo aleatorio hecho desde godot durante toda la parte del desarrollo hasta que se implementó uno mas avanzado.

El comportamiento avanzado se basa en un modelo de aprendizaje por refuerzo entrenado con \textbf{MaskablePPO} de la biblioteca \textit{SB3-Contrib} sobre un entorno \textbf{Gymnasium} que replica fielmente las reglas del videojuego de forma independiente al motor Godot. El espacio de observación se codifica mediante tres canales matriciales: el primero recoge la identidad y rango de cada pieza en el tablero, el segundo representa la transitabilidad de cada casilla, y el tercero indica la movilidad disponible para el agente en el turno actual. La máscara de acciones garantiza que el agente nunca explore movimientos inválidos durante el entrenamiento, lo que estabiliza considerablemente la convergencia. El entrenamiento emplea una estrategia de \textit{self-play}: el agente se enfrenta progresivamente a versiones anteriores de sí mismo, lo que le permite desarrollar de forma autónoma estrategias de engaño e inferencia sobre las piezas ocultas del rival.

La integración con Godot se realiza a través de un \texttt{BotController} en C\# que lanza un subproceso Python en segundo plano al iniciar la partida, carga el archivo \texttt{.zip} del modelo entrenado y se comunica con él mediante entrada y salida estándar para solicitar la acción correspondiente en cada turno. Si el subproceso no responde o el modelo no puede cargarse, el controlador activa automáticamente el modo aleatorio como fallback, garantizando la continuidad de la partida en todo momento.




